\scanpage{17}
% Three long oblique pen strokes cancel the typed passage as a whole.
\deleted{triangulée $\deleted{\unsure{a|c}}C(X,A)$ associée ainsi à $X$ détermine l'ensemble des composantes connexes de $X$, et le groupe fondamental de chacune des composantes connexes, et bien entendu la cohomologie de chacune, à coéffieients des $A$-modules tordus de type fini \ldots{} Lorsque $A=\underline Z$, il semble bien que Hurewicz (ou une variante de H. mise au point par Artin-Mazur) montre qu'un morphisme $f:X\handwritten{\to}Y$ est un isomorphisme dans la catégorie homotopique, si et seulement si le foncteur $C(Y,\underline Z)\handwritten{\to}C(\mathbb X,\underline{\mathbb Z})$ est une équivalence. On se demande alors si on ne peut pas réconstituer entièrement $X$ en termes de $C(X,\underline Z)$, du moins (si on a fait dans la définition de $C(X,\underline Z)$ les\handwritten{/}hypothèses de finitude que j'ai dites) si on se restreint à des $X$ satisfaisant certaines conditions de finitude (comme des polyèdres qui sont finis en chaque dimension). Plus précisément, serait-il vrai que tout f\unsure{a|o}ncteur exact $C(Y,\underline Z)\handwritten{\to}C(X,\underline Z)$ proviendrait à isomorphisme près d'un unique morphisme $X\handwritten{\to}Y$ ? Même si ceci était trop optimiste, les catégories $C(X,A)$ me semblent des invariants rupinants, qui doivent être des approximations très bonnes de l'objet $X$ à homotopie près, et on aimerait expliciter ce point de vue en \deleted{\unsure{explicitant|expliciter}} interprétant en \deleted{\unsure{temmes|termes}}termes de ces catégories les constructions les plus importantes familières aux homotopistes: construction des groupes d'homotopie supérieurs, plus \deleted{\unsure{génralement|généralement}}\inserted{précisément} des espaces de lacets, et en sens inverse des espaces classifiants, ou des suspensions, produits cartésiens d'espaces, \inserted{\handwritten{squelettes, \unsure{tuages|tuer|étages|images} des groupes d'homotopie supérieurs}} etc. \deleted{Notes aussi que \unsure{no|ne}} C'est là un complexe de questions manifestement très lié à celui que j'a\uline{i} abordé dans ma lettre du 26.3. dont tu viens de m'envoyer une photocopie, et je crois que quelqu'un devrait bien un jour regarder ces choses systématiquement. Notes que si tu fais sur $X$ une hypothèse de finitude comme plus haut, alors $C(X,\underline{\unsure{\mathbb Z|Z}})$ sera}
% archive sheet number: 16
